\documentclass[12pt]{article}
\usepackage[utf8]{inputenc}
\usepackage[T1]{fontenc}
\usepackage{hyperref}
\title{Sambal Language Resources for K-12 Curriculum Development}
\author{Compiled by Claude Sonnet 5.5 (claude-sonnet-5-5) for Andr\'as Kornai}
\date{}
\begin{document}
\maketitle

\section{Introduction}
Sambal (Sambali; ISO 639-3 xsb, formerly called Tina Sambal) is an Austronesian language of Zambales province, Philippines, with about 70,000 speakers (a 2000 figure). The name ``Tina'' is considered offensive by some, so ``Sambal'' is used here. Botolan Sambal is listed as a distinct variety.

\section{Primary Educational Resources}
DepEd Zambales has posted translations of several children's stories on its website, according to the search summary. I did not see a direct URL, so none is cited. No full textbook series was found.

\section{Linguistic and Reference Materials}
An \textit{English-Tina Sambal-Pilipino Dictionary} by Elgincolin and Goschnick (SIL, 1988) is cited in Wikipedia's references. Schuster and Goschnick's folktale texts with analysis appeared in \textit{Studies in Philippine Linguistics} 2(2). Overview: \url{https://en.wikipedia.org/wiki/Sambal_language}. Language capsule from UP Linguistics: \url{https://linguistics.upd.edu.ph/the-katig-collective/language-capsules/sambal} and \url{https://linguistics.upd.edu.ph/?p=9638}.

\section{Religious and Community Resources}
Most recordings are religious content hosted by Global Recordings Network, Bible.is and the Jesus Film Project (reported in search; direct URLs not seen).

\section{Digital and Audio Resources}
OLAC reportedly keeps a list of Sambal primary texts and descriptions; no direct URL was seen. Some Sambal songs are on YouTube. Online word lookup: \url{https://www.tagalog.com/dictionary/sambal}.

\section{Curriculum Development Context}
Sambal is used as a language of instruction in some schools, consistent with the Philippines' mother-tongue-based multilingual education. This search did not find a Sambal curriculum, so a next step is to contact DepEd Zambales and the SIL dictionary holders. Materials for upper grades appear to be scarce.
\end{document}